% Version: v5 (2026-10-28), the symmetrized product and the braided-symmetric derivation are recalled from braided-sym-deriv with located citations (Definition 2.1, Remark 2.4) in place of a new definition, and the symbol is BSDer as in that paper. Earlier, v4: the other paper of the series on symmetrized products is cited with a bibliography entry and a citation, in place of a web address in the text. No claim changes. Earlier, v3: sections merged (the preliminaries and the symmetrized product form one section); the standard polynomial and the Li and Wei material removed; three commutative diagrams; the companion paper is cited; one linking sentence and two self-references repaired.
\documentclass[]{interact}
\usepackage[utf8]{inputenc}

\usepackage{mathtools}
\usepackage{amsmath,amssymb,amsthm}
\usepackage{amsfonts}
\usepackage{enumerate}
\usepackage{enumitem}

\usepackage[hidelinks,pagebackref=false]{hyperref}
\usepackage[all]{hypcap}
\hypersetup{linktoc=all}

\usepackage{tikz}
\usepackage{tikz-cd}
\usepackage[yyyymmdd,hhmmss]{datetime}
\usepackage{float}

\usepackage[numbers,sort&compress]{natbib}
\bibpunct[, ]{[}{]}{,}{n}{,}{,}
\renewcommand\bibfont{\fontsize{10}{12}\selectfont}
\makeatletter
\def\NAT@def@citea{\def\@citea{\NAT@separator}}
\makeatother

\theoremstyle{plain}
\newtheorem{theorem}{Theorem}[section]
\newtheorem{lemma}[theorem]{Lemma}
\newtheorem{corollary}[theorem]{Corollary}
\newtheorem{proposition}[theorem]{Proposition}

\theoremstyle{definition}
\newtheorem{definition}[theorem]{Definition}
\newtheorem{example}[theorem]{Example}

\theoremstyle{remark}
\newtheorem{remark}[theorem]{Remark}
\newtheorem{notation}{Notation}

\allowdisplaybreaks
\newcommand{\Der}{\mathrm{Der}}
\newcommand{\Inn}{\mathrm{Inn}}
\newcommand{\Hom}{\operatorname{Hom}}
\newcommand{\JDer}{\mathrm{JDer}}
\newcommand{\BSDer}{\mathrm{BSDer}}
\newcommand{\LDer}{\mathrm{LDer}}

\newcommand{\paperversion}{5}

\begin{document}
\title{Jordan and Lie Derivations as Derivations of the Symmetrized and Antisymmetrized Products}
\author{
\name{Abderrahim Adrabi\textsuperscript{a,b}\thanks{CONTACT A. Adrabi. Email: abderrahim.adrabi@um5r.ac.ma}}
	\affil{\textsuperscript{a}Faculty of Sciences, Mohammed V University in Rabat, Morocco;\\ \textsuperscript{b}Faculty of Educational Sciences, Mohammed V University in Rabat, Morocco}
}

\history{Preprint version \paperversion, compiled \today}

\maketitle

\begin{abstract}
We view the Jordan derivations and the Lie derivations of an algebra as the derivations of two products on the algebra.
These are the symmetrized product $xy + yx$ and the antisymmetrized product $xy - yx$.
The question whether such a derivation is a derivation becomes a comparison of three spaces, $\Der(A)$, $\JDer(A)$ and $\LDer(A)$.
We show that the dimension of $\JDer(A)$ does not change when the base field is extended.
We compute the three spaces for full matrix algebras, and we apply the theorems of Zhang and Yu and of Cheung to triangular algebras, with the upper triangular $2 \times 2$ matrices as the smallest case.
In both cases a trace-type map is a Lie derivation that is not a derivation.
\end{abstract}

\begin{keywords}
Jordan derivations; Lie derivations; symmetrized product; triangular algebras; matrix algebras
\end{keywords}

\begin{amscode}
MSC 2020: 16W25; 17C50; 15A78
\end{amscode}

\section{Introduction}\label{sec:intro}

Throughout, $K$ is a field and an algebra means a unital associative $K$-algebra.
We take all maps between algebras to be $K$-linear.
We say so when a result cited below holds more generally for additive maps on rings.

A \emph{derivation} of an algebra $A$ is a linear map $D \colon A \to A$ with $D(xy) = D(x)y + xD(y)$ for all $x, y \in A$.
The notion goes back to \citet[Part I]{Jacobson1937}, who studied the derivations of an arbitrary nonassociative algebra.
\citet[\S\S2 and 4]{Hochschild1945} defined the cocycles of an associative algebra, and its $1$-cocycles are the derivations.
A \emph{Jordan derivation} of $A$ is a linear map with $D(x \circ y) = D(x) \circ y + x \circ D(y)$ for all $x, y \in A$, where $x \circ y = xy + yx$.
\citet[p.~1104]{Herstein1957} introduced this notion.
His result needs three notions for a ring $R$.
We call $R$ \emph{prime} if $aRb = 0$ implies $a = 0$ or $b = 0$ \citep[p.~1003]{Bresar1988}.
We call $R$ \emph{semiprime} if $aRa = 0$ implies $a = 0$ \citep[p.~1003]{Bresar1988}.
We call $R$ \emph{$2$-torsion free} if $2a = 0$ implies $a = 0$ \citep[p.~1003]{Bresar1988}.
So every prime ring is semiprime.
Herstein proved that every Jordan derivation of a prime ring of characteristic different from $2$ is a derivation \citep[Theorem 3.1]{Herstein1957}.
Every derivation is a Jordan derivation, and the converse is false in general \citep[p.~1003]{Bresar1988}.
\citet[Corollary 5]{Cusack1975} extended this result of Herstein to $2$-torsion free semiprime rings, and \citet[Theorem 1]{Bresar1988} gave an independent proof.

A \emph{Lie derivation} of $A$ is a linear map with $D([x,y]) = [D(x), y] + [x, D(y)]$, where $[x,y] = xy - yx$.
\citet[p.~183]{Martindale1964} studied the Lie derivations of a primitive ring of characteristic not $2$ that contains a nontrivial idempotent, and showed that each has the form $D + T$, with $D$ a derivation and $T$ a map into the centre that vanishes on commutators.
So for Lie derivations no result of Herstein type can hold.
Here is the reason.
Let $t \colon A \to Z(A)$ be a linear map into the centre with $t([x,y]) = 0$ for all $x, y$.
Every term of the Lie derivation identity for $t$ is then zero, so $t$ is a Lie derivation.
But $t$ is not a derivation as soon as $t(1) \neq 0$.

So a description of Lie derivations must allow such a term.
A Lie derivation has the \emph{standard form} if it equals $d + t$, with $d$ a derivation and $t$ a central-valued map that vanishes on commutators \citep[\S1]{Cheung2003}.
A ring is \emph{left primitive} if it has a faithful simple left module, and right primitive rings are defined in the same way \citep[Definition 11.2]{Lam2001}.
\citet[Theorem 2]{Martindale1964} proved that every Lie derivation has this form for primitive rings of characteristic not $2$ with a nontrivial idempotent.
Here an \emph{idempotent} is an element $e$ with $e^2 = e$, and it is \emph{nontrivial} if $e \neq 0$ and $e \neq 1$.
Here $d$ may take values in a larger ring.


Triangular algebras are defined in Section~\ref{sec:symmetrized}.
For them, \citet[Theorem 2.1]{ZhangYu2006} proved that every Jordan derivation of a triangular algebra is a derivation.
For Lie derivations, \citet[Theorem 11]{Cheung2003} found conditions under which every Lie derivation has the standard form.
\citet[Theorem 3.1 and Corollary 3.1]{Wang2019} treated every triangular algebra, with no condition on the module or the centre. For these algebras a Lie derivation is a sum $\delta + \tau$, where $\delta$ is a derivation into a larger triangular algebra and $\tau$ is a linear map into its centre that vanishes on commutators.


In this paper a Jordan derivation is a derivation of the symmetrized product $x \circ y$, and a Lie derivation is a derivation of the antisymmetrized product $[x,y]$. % claim-ok: states the point of view of this paper
We follow this order.
Section~\ref{sec:symmetrized} recalls the classical notions, recalls the symmetrized product, the braided-symmetric derivations and their values for the identity and the swap from the first paper, computes the dual numbers, and applies the theorem of Zhang and Yu to triangular algebras.
Section~\ref{sec:base-change} shows that the dimension of the space of Jordan derivations is stable under extension of the base field.
Section~\ref{sec:lie} treats the antisymmetrized product, full matrix algebras, and triangular algebras through the theorems of Martindale and Cheung.
We place each example right after the result it illustrates.
The twisted symmetrized product and its derivations are the subject of the first paper in our series on symmetrized products~\cite{AdrabiBennis2026a}, from which the definition and the dictionary of Section~\ref{sec:symmetrized} are taken.

\section{Preliminaries and the symmetrized product}\label{sec:symmetrized}

This section recalls the classical notions we need, then the symmetrized product and the braided-symmetric derivations of the first paper, with the result of that paper for the identity and the swap.


From now on, $A$ is a $K$-algebra with unit $1$.
We write $[x,y] = xy - yx$ for the commutator, $x \circ y = xy + yx$ for the Jordan product, and $Z(A)$ for the centre.
The centre is the set of elements $z$ with $zx = xz$ for all $x \in A$.
We call an element $x$ \emph{nilpotent} if $x^n = 0$ for some $n \geq 1$.
A commutative algebra is \emph{reduced} if its only nilpotent element is zero \citep[1.4.12]{ChambertLoir2021}.

A derivation of an algebra is as defined by \citet[Part I]{Jacobson1937}: a linear map $D$ with $D(xy) = D(x)y + xD(y)$.
Every derivation satisfies $D(1) = 0$, since $D(1) = D(1 \cdot 1) = 2D(1)$.
For $g \in A$, the map $\operatorname{ad}_g(x) = [g,x]$ is an \emph{inner derivation}, called a principal derivation by \citet[\S9.2]{Weibel1994}.
We write $\Der(A)$ and $\Inn(A)$ for these two spaces.

We add a word on the characteristic for Jordan derivations.
A \emph{Jordan derivation} in the sense of \citet[p.~1104]{Herstein1957} is an additive map with $D(x^2) = D(x)x + xD(x)$.
We write $\JDer(A)$ for the space of linear maps with $D(x \circ y) = D(x) \circ y + x \circ D(y)$ for all $x, y$.
Replacing $x$ by $x + y$ gives the condition $D(x \circ y) = D(x) \circ y + x \circ D(y)$ in every characteristic.
Setting $y = x$ in this condition gives $2D(x^2) = 2(D(x)x + xD(x))$.
So the two conditions agree when $2$ is invertible in $K$, and we use the $\circ$ form throughout.
When $\operatorname{char} K = 2$ and $A$ is commutative, $x \circ y = 0$, and every linear map satisfies the $\circ$ condition.

We need two more kinds.
A \emph{Lie derivation} satisfies $D([x,y]) = [D(x), y] + [x, D(y)]$ \citep[p.~183]{Martindale1964}, and we write $\LDer(A)$ for this space.
When $A$ is commutative, every bracket is zero, so every linear map is a Lie derivation.

We now compare these spaces.
When $2$ is invertible, $xy = \tfrac12 (x \circ y) + \tfrac12 [x,y]$.
A derivation respects $xy$, so it respects both halves, and $\Der(A) \subseteq \JDer(A) \cap \LDer(A)$.
Both inclusions can be strict, as Examples~\ref{ex:dual-numbers} and~\ref{ex:matrix} show.
Our examples include triangular algebras.
An \emph{$(A,B)$-bimodule} is a left $A$-module and right $B$-module $M$ with $(am)b = a(mb)$ for all $a \in A$, $m \in M$ and $b \in B$ \citep[Example 1.14]{Lam2001}.
We call it \emph{unital} if $1m = m = m1$ for all $m \in M$.

Let $R$ be a commutative ring, let $A$ and $B$ be unital $R$-algebras, and let $M$ be a unital $(A,B)$-bimodule.
The \emph{triangular algebra} $T = \mathrm{Tri}(A,M,B)$ is the set of matrices $\left(\begin{smallmatrix} a & m \\ 0 & b \end{smallmatrix}\right)$ with the matrix sum and product \citep[\S2]{Cheung2003}.
The module $M$ is \emph{faithful} as a left $A$-module if $aM = 0$ implies $a = 0$, and as a right $B$-module if $Mb = 0$ implies $b = 0$.
The idempotents $e_1 = \left(\begin{smallmatrix} 1 & 0 \\ 0 & 0 \end{smallmatrix}\right)$ and $e_2 = \left(\begin{smallmatrix} 0 & 0 \\ 0 & 1 \end{smallmatrix}\right)$ satisfy $e_2 T e_1 = 0$.
So $T$ is not prime, since $e_1 \neq 0$ and $e_2 \neq 0$, and the prime-ring results above do not apply to triangular algebras.



Let $A$ be a $K$-algebra with product $\mu$.
Let $\sigma \colon A \otimes A \to A \otimes A$ be a $K$-linear map with $\sigma^2 = \mathrm{id}$, for instance the swap $\tau(x \otimes y) = y \otimes x$.
It gives a symmetrized product \cite[Section 2]{AdrabiBennis2026a}
\[
x \cdot_\sigma y = \mu(x \otimes y) + \mu(\sigma(x \otimes y)).
\]
For $\sigma = \tau$ this is the Jordan product $x \circ y$.

\begin{definition}[{\cite[Definition 2.1]{AdrabiBennis2026a}}]\label{def:sigma-derivation}
Let $A$ be a $K$-algebra and let $\sigma$ be a $K$-linear involution of $A \otimes A$.
A \emph{braided-symmetric derivation} for $\sigma$ is a $K$-linear map $D \colon A \to A$ with
\begin{equation}\label{eq:sigma-leibniz}
D(x \cdot_\sigma y) = D(x) \cdot_\sigma y + x \cdot_\sigma D(y) \quad \text{for all } x, y \in A.
\end{equation}
We write $\BSDer_\sigma(A)$ for the set of braided-symmetric derivations for $\sigma$.
This set is a $K$-subspace of $\mathrm{End}_K(A)$, since the condition is linear in $D$.
For an automorphism $\alpha$ of $A$, we write $\BSDer_\sigma(A)^\alpha$ for the subset of maps $D$ with $D\alpha = \alpha D$.
\end{definition}

Identity~\eqref{eq:sigma-leibniz} says that the following square commutes.
\[
\begin{tikzcd}[column sep=large]
A \otimes A \arrow[r, "\mu_\sigma"] \arrow[d, "D \otimes 1 + 1 \otimes D"'] & A \arrow[d, "D"] \\
A \otimes A \arrow[r, "\mu_\sigma"'] & A
\end{tikzcd}
\]

The next remark, recalled from the first paper, says what Definition~\ref{def:sigma-derivation} gives for the two involutions that use no structure of $A$.

\begin{remark}[{\cite[Remark 2.4]{AdrabiBennis2026a}}]\label{rem:sigma-derivation-classification}
Let $A$ be a $K$-algebra and let $\sigma$ be an involution of $A \otimes A$.
Identity~\eqref{eq:sigma-leibniz} for $\sigma$ is the Leibniz rule for the product $\cdot_\sigma$.
So $\BSDer_\sigma(A)$ is the set of derivations of the nonassociative algebra $(A, \cdot_\sigma)$, in the sense of \citet[Part I]{Jacobson1937}.

For the identity, $x \cdot_{\mathrm{id}} y = 2xy$, so the identity reads
\[
2D(xy) = 2\bigl(D(x)y + xD(y)\bigr).
\]
If $\operatorname{char} K \neq 2$, we divide by $2$ and get $\BSDer_{\mathrm{id}}(A) = \Der(A)$.
If $\operatorname{char} K = 2$, then $x \cdot_{\mathrm{id}} y = 0$, so $\BSDer_{\mathrm{id}}(A)$ contains every linear map.
So it is larger than $\Der(A)$.
For example, take the linear map with $D(1) = 1$ and $D = 0$ on a complement of $K1$.
Then $D$ is not a derivation, since a derivation satisfies $D(1) = D(1 \cdot 1) = 2D(1) = 0$ in characteristic $2$.

For the swap $\tau$ we get $x \cdot_\tau y = xy + yx = x \circ y$.
So the defining identity is the $\circ$ form of the Jordan identity recalled above.
Hence $\BSDer_\tau(A) = \JDer(A)$ in every characteristic.

Every derivation is a Jordan derivation \citep[p.~1003]{Bresar1988}.
So in every characteristic
\[
\Der(A) \subseteq \BSDer_\tau(A) = \JDer(A).
\]
When $\operatorname{char} K \neq 2$, this gives $\BSDer_{\mathrm{id}}(A) \subseteq \BSDer_\tau(A)$.
This is an inclusion between the spaces of two involutions on one algebra. % claim-ok: names the inclusion displayed just above
\end{remark}

The smallest nonreduced algebra shows both inclusions recalled above at work.

\begin{example}[Dual numbers]\label{ex:dual-numbers}
Let $A = K[\varepsilon]/(\varepsilon^2)$, with basis $\{1, \varepsilon\}$.
A linear map $D$ is a derivation exactly when the Leibniz rule holds on the four pairs of basis elements.
Since $D(1) = 0$, the pairs $(1,\varepsilon)$ and $(\varepsilon,1)$ hold.
Write $D(\varepsilon) = \beta + c\varepsilon$ with $\beta, c \in K$.
The pair $(\varepsilon,\varepsilon)$ gives $0 = D(\varepsilon^2) = 2\varepsilon D(\varepsilon) = 2\beta\varepsilon$.
If $\operatorname{char} K \neq 2$, then $\beta = 0$, and every derivation has the form $D(a + b\varepsilon) = bc\varepsilon$, so $\dim_K \Der(A) = 1$.
If $\operatorname{char} K = 2$, then $\beta$ is free, and $\dim_K \Der(A) = 2$.
Since $A$ is commutative, every commutator is zero, so $\LDer(A) = \mathrm{End}_K(A)$, which has dimension $2 \cdot 2 = 4$.
Each derivation lies in $\BSDer_\tau(A)$ by Remark~\ref{rem:sigma-derivation-classification}.
If $\operatorname{char} K = 2$, then $x \circ y = 2xy = 0$ for all $x, y$, so $\JDer(A) = \mathrm{End}_K(A)$ has dimension $4$, which exceeds $\dim_K \Der(A) = 2$.
\end{example}

The remark above and the following theorem on triangular algebras give a corollary. % claim-ok: signpost to the corollary proved right after

\begin{theorem}[{\citep[Theorem 2.1]{ZhangYu2006}}]\label{thm:zhang-yu}
Let $R$ be a $2$-torsion free commutative ring, let $A$ and $B$ be unital $R$-algebras, and let $M$ be a unital $(A,B)$-bimodule that is faithful as a left $A$-module and as a right $B$-module. Then every Jordan derivation $D \colon \mathrm{Tri}(A,M,B) \to \mathrm{Tri}(A,M,B)$ is a derivation.
\end{theorem}

In the next corollary, maps on $T = \mathrm{Tri}(A,M,B)$ are $R$-linear and the $\sigma$-symmetric identity is read over $R$.
It follows from Remark~\ref{rem:sigma-derivation-classification} and Theorem~\ref{thm:zhang-yu}.

\begin{corollary}\label{cor:triangular-jordan} % ours: follows from the remark and the theorem of Zhang and Yu stated above
Let $R$ be a $2$-torsion free commutative ring, let $A$ and $B$ be unital $R$-algebras, and let $M$ be a unital $(A,B)$-bimodule that is faithful as a left $A$-module and as a right $B$-module.
Suppose also that $\mathrm{Tri}(A,M,B)$ is $2$-torsion free.
Then $\BSDer_\tau(\mathrm{Tri}(A,M,B)) = \Der(\mathrm{Tri}(A,M,B))$.
\end{corollary}

\section{Extension of the base field}\label{sec:base-change}

Proposition~\ref{prop:base-change} below treats extension of the base field.
For a field extension $L$ of $K$ and a $K$-vector space $V$, write $V_L = L \otimes_K V$.
This is an $L$-vector space. % claim-ok: the tensor product of a $K$-vector space with $L$ over $K$ is an $L$-vector space by construction
For a $K$-algebra $A$, the space $A_L$ is an $L$-algebra with $(c \otimes x)(c' \otimes y) = cc' \otimes xy$ \citep[Chapter 2, \S5, p.~74]{Lam2001}.

We now define the involution $\sigma_L$ of $A_L \otimes_L A_L$.
It is an $L$-linear map of $A_L \otimes_L A_L$ into itself with $\sigma_L^2 = \mathrm{id}$.

Define
\[
\Psi \colon A_L \otimes_L A_L \to L \otimes_K (A \otimes_K A), \qquad
(c \otimes x) \otimes_L (c' \otimes y) \longmapsto cc' \otimes (x \otimes y).
\]
Since this assignment is $L$-bilinear, it defines an $L$-linear map.
It is an isomorphism, because its inverse sends $c \otimes (x \otimes y)$ to $(c \otimes x) \otimes_L (1 \otimes y)$.

The map $\mathrm{id}_L \otimes \sigma$ is an $L$-linear map of $L \otimes_K (A \otimes_K A)$ into itself, and $(\mathrm{id}_L \otimes \sigma)^2 = \mathrm{id}_L \otimes \sigma^2 = \mathrm{id}$.
We define $\sigma_L$ as the conjugate of $\mathrm{id}_L \otimes \sigma$ by $\Psi$:
\[
\sigma_L = \Psi^{-1} \circ (\mathrm{id}_L \otimes \sigma) \circ \Psi.
\]
For $c, c' \in L$ and $x, y \in A$ it reads
\[
\sigma_L\bigl((c \otimes x) \otimes_L (c' \otimes y)\bigr) = \Psi^{-1}\bigl(cc' \otimes \sigma(x \otimes y)\bigr).
\]

Since $\Psi \Psi^{-1} = \mathrm{id}$, we get
\[
\sigma_L^2 = \Psi^{-1} (\mathrm{id}_L \otimes \sigma)(\mathrm{id}_L \otimes \sigma) \Psi = \Psi^{-1} (\mathrm{id}_L \otimes \sigma^2) \Psi = \mathrm{id}.
\]
So $\sigma_L$ is an involution of $A_L \otimes_L A_L$.

The product $\mu_L$ of $A_L$ satisfies $\mu_L = (\mathrm{id}_L \otimes \mu) \circ \Psi$.
So $\Psi$ carries $\mu_L$ and $\sigma_L$ to $\mathrm{id}_L \otimes \mu$ and $\mathrm{id}_L \otimes \sigma$.
Hence $(1 \otimes x) \cdot_{\sigma_L} (1 \otimes y) = 1 \otimes (x \cdot_\sigma y)$ for $x, y \in A$.
Thus the product $\cdot_{\sigma_L}$ on $A_L$ extends the product $\cdot_\sigma$ on $A$.

Proposition~\ref{prop:base-change} below says that extending the base field does not change the dimension of the space of braided-symmetric derivations. The proof compares spaces of linear maps over $K$ and over $L$, and uses results from the recent literature, which we recall. % claim-ok: signpost for the cited results below
For a $K$-vector space $E$ we write $E_L = L \otimes_K E$, and for a $K$-linear map $f$ we write $f_L = \mathrm{id}_L \otimes f$.

Let $L$ be an extension field of $K$, and let $(c_i)_{i \in I}$ be a basis of $L$ over $K$.
For a $K$-linear map $f \colon E \to F$, every element of $E_L$ is a unique sum $\sum_i c_i \otimes e_i$ with $e_i \in E$, and $f_L(\sum_i c_i \otimes e_i) = \sum_i c_i \otimes f(e_i)$.
So $f_L$ kills the element exactly when every $f(e_i)$ is zero, and $\ker(f_L) = (\ker f)_L$.
% claim-ok: direct computation on a basis of L over K

The next two results compare linear maps over $K$ and over $L$, which is the comparison that Proposition~\ref{prop:base-change} needs to transfer the equation $\Phi(D) = 0$ from $A$ to $A_L$. % claim-ok: signpost for the lemmas below

\begin{lemma}[{\cite[Lemma 4.86]{Rotman2009}}]\label{lem:rotman-486}
Let $B$ be a flat $R$-module, and let $N$ be a finitely presented $R$-module.
For every $R$-module $M$, there is a natural isomorphism
\[
\psi \colon B \otimes_R \mathrm{Hom}_R(N, M) \to \mathrm{Hom}_R(N, M \otimes_R B),
\]
given by $b \otimes g \mapsto g_b$, where $g_b(n) = g(n) \otimes b$ for all $n \in N$.
\end{lemma}

The next lemma identifies homomorphisms over $R$ with homomorphisms over $A$ after extension of scalars. % claim-ok: signpost for the lemma stated right after

\begin{lemma}[{\cite[Lemma 4.85]{Rotman2009}}]\label{lem:rotman-485}
Let $A$ be an $R$-algebra, and let $N$ be a finitely presented $R$-module.
For every $A$-module $M$, there is a natural isomorphism
\[
\theta \colon \mathrm{Hom}_R(N, M) \to \mathrm{Hom}_A(N \otimes_R A, M),
\]
given by $\theta \colon f \mapsto \tilde f$, where $\tilde f(n \otimes 1) = f(n)$ for all $n \in N$.
\end{lemma}

Let $V$ be a finite-dimensional $K$-vector space and let $W$ be a $K$-vector space.
Then $V$ is finitely presented over $K$, by a finite basis. % claim-ok: a finite basis gives a presentation
Lemma~\ref{lem:rotman-486} with $R = K$, $B = L$, $N = V$ and $M = W$ gives $L \otimes_K \mathrm{Hom}_K(V, W) \cong \mathrm{Hom}_K(V, W_L)$.
Lemma~\ref{lem:rotman-485} with $R = K$, $A = L$, $N = V$ and $M = W_L$ gives $\mathrm{Hom}_K(V, W_L) \cong \mathrm{Hom}_L(V_L, W_L)$.
The composite sends $c \otimes u$ to $c\,(\mathrm{id}_L \otimes u)$. % claim-ok: the two maps of the lemmas written on the generators
So $L \otimes_K \mathrm{Hom}_K(V, W) \to \mathrm{Hom}_L(V_L, W_L)$ is bijective, which identifies $L \otimes_K \mathrm{End}_K(A)$ with $\mathrm{End}_L(A_L)$ in Proposition~\ref{prop:base-change}. % claim-ok: conclusion of the two lines above

For a $K$-algebra $A$ with involution $\sigma$ of $A \otimes A$, let $m_\sigma \colon A \otimes A \to A$ be the product $x \otimes y \mapsto x \cdot_\sigma y$.
Define the $K$-linear map
\[
\Phi \colon \mathrm{End}_K(A) \to \mathrm{Hom}_K(A \otimes A, A), \qquad
\Phi(D) = D \circ m_\sigma - m_\sigma \circ (D \otimes \mathrm{id} + \mathrm{id} \otimes D).
\]
So $\Phi(D)$ is the difference of the two sides of identity~\eqref{eq:sigma-leibniz}, and $\BSDer_\sigma(A) = \ker \Phi$.
Write $\Phi_L$ for the same map for $A_L$ and $\sigma_L$, so $\BSDer_{\sigma_L}(A_L) = \ker \Phi_L$.

\begin{proposition}\label{prop:base-change}
Let $K$ be a field, let $A$ be a finite-dimensional $K$-algebra, and let $\sigma$ be an involution of $A \otimes A$.
Let $L$ be a field extension of $K$.
For $D \in \mathrm{End}_K(A)$ and $c \in L$, let $c \otimes D$ act on $A_L$ by $(c \otimes D)(c' \otimes x) = cc' \otimes D(x)$ for $c' \in L$ and $x \in A$.
This rule defines isomorphisms of $L$-vector spaces
\[
\beta \colon L \otimes_K \mathrm{End}_K(A) \to \mathrm{End}_L(A_L), \qquad
\gamma \colon L \otimes_K \mathrm{Hom}_K(A \otimes A, A) \to \mathrm{Hom}_L(A_L \otimes_L A_L, A_L)
\]
such that the following square commutes.
\[
\begin{tikzcd}[column sep=huge]
L \otimes_K \mathrm{End}_K(A) \arrow[r, "\mathrm{id}_L \otimes \Phi"] \arrow[d, "\beta"'] & L \otimes_K \mathrm{Hom}_K(A \otimes A, A) \arrow[d, "\gamma"] \\
\mathrm{End}_L(A_L) \arrow[r, "\Phi_L"'] & \mathrm{Hom}_L(A_L \otimes_L A_L, A_L)
\end{tikzcd}
\]
So $\beta$ and $\gamma$ carry the exact sequence
\[
0 \to L \otimes_K \BSDer_\sigma(A) \to L \otimes_K \mathrm{End}_K(A) \xrightarrow{\ \mathrm{id}_L \otimes \Phi\ } L \otimes_K \mathrm{Hom}_K(A \otimes A, A)
\]
isomorphically onto the exact sequence
\[
0 \to \BSDer_{\sigma_L}(A_L) \to \mathrm{End}_L(A_L) \xrightarrow{\ \Phi_L\ } \mathrm{Hom}_L(A_L \otimes_L A_L, A_L).
\]
In particular $L \otimes_K \BSDer_\sigma(A) \cong \BSDer_{\sigma_L}(A_L)$ and $\dim_L \BSDer_{\sigma_L}(A_L) = \dim_K \BSDer_\sigma(A)$.
\end{proposition}

\begin{proof}
First, the spaces $A$ and $A \otimes_K A$ are finite-dimensional.
So $\beta$ is an isomorphism by the bijection above with $V = W = A$.
Also $\gamma$ is an isomorphism by the bijection above with $V = A \otimes_K A$ and $W = A$, since $(A \otimes_K A)_L = A_L \otimes_L A_L$.

Second, we check $\Phi_L \circ \beta = \gamma \circ (\mathrm{id}_L \otimes \Phi)$.
Both sides are $L$-linear.
So it is enough to evaluate them on $c \otimes D$ and on $(c' \otimes x) \otimes_L (c'' \otimes y)$, for $c, c', c'' \in L$, $D \in \mathrm{End}_K(A)$ and $x, y \in A$.
By the construction of $\sigma_L$, we have $m_{\sigma_L}\bigl((c' \otimes x) \otimes_L (c'' \otimes y)\bigr) = c'c'' \otimes m_\sigma(x \otimes y)$.
Hence both sides give
\[
cc'c'' \otimes \Bigl(D\bigl(m_\sigma(x \otimes y)\bigr) - m_\sigma\bigl(D(x) \otimes y + x \otimes D(y)\bigr)\Bigr) = cc'c'' \otimes \Phi(D)(x \otimes y).
\]

Third, $\ker \Phi = \BSDer_\sigma(A)$.
Since $L \otimes_K -$ is exact, $\ker(\mathrm{id}_L \otimes \Phi) = L \otimes_K \BSDer_\sigma(A)$.
By the second step, $\beta$ and $\gamma$ turn $\mathrm{id}_L \otimes \Phi$ into $\Phi_L$.
So $\beta$ maps $L \otimes_K \BSDer_\sigma(A)$ onto $\ker \Phi_L = \BSDer_{\sigma_L}(A_L)$.
\end{proof}


The next remark shows that the finite-dimensional hypothesis of Proposition~\ref{prop:base-change} is needed.

\begin{remark}\label{rem:finite-needed}
The finite-dimensional hypothesis in Proposition~\ref{prop:base-change} cannot be dropped.
Take $\sigma = -\mathrm{id}$.
Then $\sigma^2 = \mathrm{id}$, so $\sigma$ is an involution of $A \otimes A$, and $x \cdot_\sigma y = 0$ for all $x, y$.
So every $K$-linear map is a braided-symmetric derivation, both for $A$ and for $A_L$.

Let $A = K[x]$ and $L = K(t)$, where $t$ is transcendental over $K$.
Define $D' \in \mathrm{End}_L(A_L)$ by $D'(x^i) = t^i x^i$ for all $i \geq 0$.
Suppose $D' = \sum_{r=1}^s c_r \otimes D_r$ for some $s \geq 1$, with $c_r \in L$ and $D_r \in \mathrm{End}_K(A)$.
Since $D_r$ is $K$-linear and $A = K[x]$, the coefficient of $x^i$ in $D_r(x^i)$ lies in $K$.
So $t^i$ lies in the $K$-span of $c_1, \dots, c_s$, for every $i$.
Since $t$ is transcendental over $K$, the powers $1, t, t^2, \dots$ are linearly independent over $K$, so no finite span contains them all.
Hence $D'$ is not in the image, and the map of Proposition~\ref{prop:base-change} is not onto.
\end{remark}

The next corollary applies the proposition to the space of Jordan derivations. % claim-ok: signpost for the result stated right after

\begin{corollary}\label{cor:jder-base}
Let $K$ be a field, let $A$ be a finite-dimensional $K$-algebra, and let $L$ be a field extension of $K$.
Then $\dim_L \JDer(A_L) = \dim_K \JDer(A)$.
\end{corollary}

\begin{proof}
For the swap $\tau$ of $A \otimes A$, the involution $\tau_L$ is the swap of $A_L \otimes_L A_L$, since $\tau_L((c \otimes x) \otimes_L (c' \otimes y)) = \Psi^{-1}(cc' \otimes (y \otimes x)) = (c' \otimes y) \otimes_L (c \otimes x)$.
By Remark~\ref{rem:sigma-derivation-classification}, $\BSDer_\tau(A) = \JDer(A)$ and $\BSDer_{\tau_L}(A_L) = \JDer(A_L)$.
Proposition~\ref{prop:base-change} gives the equality of dimensions.
\end{proof}

\section{The antisymmetrized product and Lie derivations}\label{sec:lie}

This section replaces the sum $\mu + \mu \circ \sigma$ by the difference $\mu - \mu \circ \sigma$.
For the swap the difference is the commutator, and its derivations are the Lie derivations.

\begin{definition}\label{def:anti-sigma-derivation} % ours: the antisymmetrized counterpart of the notion of this paper
The \emph{$\sigma$-antisymmetrized product} on a $K$-algebra $A$ is $x \cdot_\sigma^- y = \mu(x \otimes y) - \mu(\sigma(x \otimes y))$, where $\sigma$ is an involution of $A \otimes A$.
A \emph{$\sigma$-antisymmetric derivation} is a $K$-linear map $D \colon A \to A$ with $D(x \cdot_\sigma^- y) = D(x) \cdot_\sigma^- y + x \cdot_\sigma^- D(y)$.
\end{definition}

We now identify these derivations for the swap.
It also gives the easy half of the theorems of Martindale and Bre\v{s}ar.

\begin{proposition}\label{prop:lie-classification}
Let $A$ be an associative $K$-algebra.
\begin{enumerate}[label=(\roman*)]
\item The $\sigma$-antisymmetric derivations for $\tau$ in the sense of Definition~\ref{def:anti-sigma-derivation} are exactly the Lie derivations of $A$.
\item $\Der(A) \subseteq \LDer(A)$.
\item Let $d \in \Der(A)$, and let $t \colon A \to Z(A)$ be linear with $t([x,y]) = 0$ for all $x, y$. Then $d + t \in \LDer(A)$.
\item In (iii), if $t(1) \neq 0$, then $d + t \notin \Der(A)$.
\end{enumerate}
\end{proposition}

\begin{proof}
We check the four parts in order.
(i) For the swap, $x \cdot_\tau^- y = [x,y]$, so the two identities are the same.
(ii) For $D \in \Der(A)$, $D([x,y]) = D(x)y + xD(y) - D(y)x - yD(x) = [D(x), y] + [x, D(y)]$.
(iii) By (ii), $d([x,y]) = [d(x),y] + [x,d(y)]$, and $t([x,y]) = 0$.
Since $t(x)$ and $t(y)$ are central, $[t(x), y] = 0 = [x, t(y)]$.
Adding gives the Lie derivation identity for $d + t$.
(iv) Every derivation vanishes at $1$, as recalled in Section~\ref{sec:symmetrized}, while $(d+t)(1) = t(1) \neq 0$.
\end{proof}

So inner derivations lie in both spaces.
Each $\operatorname{ad}_g$ is a Lie derivation by Proposition~\ref{prop:lie-classification}(ii), and a Jordan derivation by Remark~\ref{rem:sigma-derivation-classification}.
Full matrix algebras show that part (iii) gives maps outside $\Der(A)$.
We use one result on Jordan derivations of matrix rings.
It is stated for a unital associative ring and a $2$-torsion free bimodule \citep[\S2]{Alizadeh2009}.

\begin{theorem}[{\citep[Theorem 3.1]{Alizadeh2009}}]\label{thm:alizadeh}
Let $A$ be a unital associative ring, let $M$ be a $2$-torsion free $A$-bimodule, and let $n \geq 2$. Then every Jordan derivation $D \colon M_n(A) \to M_n(M)$ is a derivation.
\end{theorem}

We also use a description of the first Hochschild cohomology.
It gives $HH^1$ of the matrix algebra in Example~\ref{ex:matrix}.
Let $R$ be a $k$-algebra and $M$ an $R$-bimodule.
\citet[\S9.2]{Weibel1994} writes $\mathrm{Der}_k(R,M)$ for the $k$-derivations.
Weibel writes $\mathrm{PDer}_k(R,M)$ for the principal ones $f_m(r) = rm - mr$.

\begin{lemma}[{\citep[Lemma 9.2.1]{Weibel1994}}]\label{lem:weibel}
Let $R$ be a $k$-algebra and $M$ an $R$-bimodule. Then $HH^1(R, M) \cong \mathrm{Der}_k(R,M)/\mathrm{PDer}_k(R,M)$.
\end{lemma}

A nonzero ring $R$ is \emph{simple} if $(0)$ and $R$ are the only ideals in $R$, and a commutative ring is simple if and only if it is a field \citep[\S1]{Lam2001}.

The next theorem shows that $M_n(K)$ is simple.
Example~\ref{ex:matrix} needs this twice: to apply the theorem of Martindale, and to see that $M_n(K)$ is separable.

\begin{theorem}[{\citep[Theorem 3.1]{Lam2001}}]\label{thm:lam-matrix}
Let $R$ be a ring and $M_n(R)$ be the ring of $n \times n$ matrices over $R$. Then any ideal $I$ of $M_n(R)$ has the form $M_n(\mathfrak{m})$ for a uniquely determined ideal $\mathfrak{m}$ of $R$. In particular, if $R$ is a simple ring, so is $M_n(R)$.
\end{theorem}

The next theorem of Martindale describes the Lie derivations of primitive rings with a nontrivial idempotent. % claim-ok: signpost for the result stated right after

\begin{theorem}[{\citep[Theorem 2]{Martindale1964}}]\label{thm:martindale}
Let $D$ be a Lie derivation of a primitive ring $R$ into itself, where $R$ contains a nontrivial idempotent and the characteristic of $R$ is not $2$. Then $D$ is of the form $d + t$, where $d$ is an ordinary derivation of $R$ into a primitive ring $\hat{R}$ containing $R$ and $t$ is an additive mapping of $R$ into the center of $\hat{R}$ that maps commutators into zero.
\end{theorem}

Martindale remarks after Theorem~\ref{thm:martindale} that, if $R$ is simple, then $d$ maps $R$ into itself and $t$ maps $R$ into the center of $R$ \citep[p.~187]{Martindale1964}.

\begin{definition}[{\citep[Section 9.2.1]{Weibel1994}}]\label{def:separable}
A finite-dimensional semisimple algebra $R$ over a field $k$ is called \emph{separable} if for every extension field $\ell$ of $k$ the $\ell$-algebra $R_\ell = R \otimes_k \ell$ is semisimple.
\end{definition}

The next theorem of Weibel characterizes separable algebras by the vanishing of Hochschild cohomology. % claim-ok: signpost for the result stated right after

\begin{theorem}[{\citep[Theorem 9.2.11]{Weibel1994}}]\label{thm:hochschild}
Let $R$ be an algebra over a field $k$.
The following are equivalent:
\begin{enumerate}
\item $R$ is a finite-dimensional separable $k$-algebra.
\item $R$ is projective as a left $R^e$-module.
\item $H_*(R, M) = 0$ for all $* \neq 0$ and all bimodules $M$.
\item $H^*(R, M) = 0$ for all $* \neq 0$ and all bimodules $M$.
\end{enumerate}
Here $R^e = R \otimes_k R^{\mathrm{op}}$, and $H_*$ and $H^*$ are Hochschild homology and cohomology.
\end{theorem}

We now apply Theorems~\ref{thm:alizadeh}, \ref{thm:lam-matrix}, \ref{thm:martindale} and~\ref{thm:hochschild} to $M_n(K)$.

\begin{example}[Full matrix algebras]\label{ex:matrix}
Let $A = M_n(K)$ with $n \geq 2$, and let $e_{ij}$ be the matrix unit with entry $1$ in position $(i,j)$ and $0$ elsewhere.
If $\operatorname{char} K \neq 2$, every Jordan derivation of $M_n(K)$ is a derivation by Theorem~\ref{thm:alizadeh} applied to the ring $K$ and the $K$-bimodule $K$, since $K$ is then $2$-torsion free.
So $\BSDer_\tau(A) = \JDer(A) = \Der(A)$ in that case.

Let $\theta(x) = \operatorname{tr}(x) I$.
Its values are central, and $\theta([x,y]) = (\operatorname{tr}(xy) - \operatorname{tr}(yx)) I = 0$.
By Proposition~\ref{prop:lie-classification}(iii), $\operatorname{ad}_g + c\,\theta$ is a Lie derivation for every $g \in A$ and $c \in K$.
Since $\theta(e_{11}) = I$ while $\theta(e_{11})e_{11} + e_{11}\theta(e_{11}) = 2e_{11}$ and $\theta(e_{11}^2) = \theta(e_{11})$, the map $\theta$ is not a derivation in any characteristic.

Now let $n \geq 2$ and $\operatorname{char} K \neq 2$.
The ring $M_n(K)$ is simple by Theorem~\ref{thm:lam-matrix}, since a field is simple.
It is primitive, because $K^n$ is a faithful simple left module over it, and $e_{11}$ is a nontrivial idempotent.
By Theorem~\ref{thm:martindale} and the remark after it, every Lie derivation $D$ of $M_n(K)$ has the form $D = d + t$, where $d$ is an additive derivation of $M_n(K)$ into itself and $t$ is an additive map into $K \cdot I$ that vanishes on commutators.

The maps $d$ and $t$ are $K$-linear.
Indeed, for $\lambda \in K$ we have $d(\lambda x) = d(\lambda I)x + \lambda d(x)$, and $d(\lambda I)$ is central because $[d(\lambda I), x] = d([\lambda I, x]) - [\lambda I, d(x)] = 0$.
So $d(\lambda I) = \nu(\lambda) I$ for some $\nu(\lambda) \in K$.
Since $D$ is $K$-linear, $\nu(\lambda)x = \lambda t(x) - t(\lambda x)$ lies in $K \cdot I$ for every $x$.
Taking $x = e_{11}$ gives $\nu(\lambda) = 0$.

Next, $M_n(K)$ is finite-dimensional, and $M_n(K) \otimes_K F = M_n(F)$ is simple for every extension field $F$ by Theorem~\ref{thm:lam-matrix}, hence semisimple. % claim-ok: a finite-dimensional simple algebra is semisimple
So $M_n(K)$ is separable by Definition~\ref{def:separable}.
By Theorem~\ref{thm:hochschild}, $H^1(M_n(K), M_n(K)) = 0$, so by Lemma~\ref{lem:weibel} every $K$-linear derivation is inner, and $d$ is inner.
Hence $\Der(A) = \Inn(A)$, and every Lie derivation of $A$ has the form $\operatorname{ad}_g + t$ with $t \colon A \to K \cdot I$ linear and $t([A,A]) = 0$.
Such a map $t$ is the same as a $K$-linear map $A/[A,A] \to K$, and $A/[A,A] = HH_0(A,A)$ \citep[\S9.1]{Weibel1994}.
A map $t$ of this kind that is also a derivation is zero.
Indeed $t(e_{ii}) = t(e_{ii}^2) = 2t(e_{ii})e_{ii}$ forces $t(e_{ii}) = 0$ when $n \geq 2$, and $t(e_{ij}) = 0$ for $i \neq j$ because $e_{ij} = [e_{ii}, e_{ij}]$.
So $t$ is determined by $\operatorname{ad}_g + t$: if $\operatorname{ad}_g + t = \operatorname{ad}_h + t'$, then $t - t' = \operatorname{ad}_{h-g}$ is such a derivation, hence zero.
Let $\omega \colon \LDer(A) \to \mathrm{Hom}_K(HH_0(A,A), K)$ send $\operatorname{ad}_g + t$ to $t$.
The map $\omega$ is surjective, since each $t$ is a Lie derivation by Proposition~\ref{prop:lie-classification}(iii), and its kernel is $\Inn(A) = \Der(A)$.
So the sequence
\[
\begin{tikzcd}[column sep=large]
0 \arrow[r] & \Der(A) \arrow[r] & \LDer(A) \arrow[r, "\omega"] & \mathrm{Hom}_K\bigl(HH_0(A,A), K\bigr) \arrow[r] & 0
\end{tikzcd}
\]
is exact.
The map $\theta$ above is one element of the right-hand space.
\end{example}

The swap gives $\JDer(A)$ through $\mu + \mu \circ \tau$ and $\LDer(A)$ through $\mu - \mu \circ \tau$.
So $\mu$ is half the sum of these two products when $2$ is invertible.
So when $2$ is invertible, a map that is both a Jordan and a Lie derivation is a derivation.
On $M_n(K)$ with $\operatorname{char} K \neq 2$ this gives $\JDer(A) \cap \LDer(A) = \Der(A)$ even though $\LDer(A)$ is larger.

For triangular algebras the Lie side has a central term.

A Lie derivation is \emph{proper} when it is the sum of a derivation and a linear map into the centre \citep[\S1]{Cheung2003}.
The extra map vanishes on commutators, so a proper Lie derivation has the standard form of Section~\ref{sec:intro}.

The next theorem uses the algebra $\mathfrak{A} = \mathrm{Tri}(A,M,B)$, and the projections $\pi_A$ and $\pi_B$ onto the diagonal blocks.

The algebra $W(X)$ is the smallest subalgebra of $X$ that satisfies seven closure conditions of \citet[Proposition 10]{Cheung2003}.
We list them for a subalgebra $V$ of $X$.
Here $F$ is a polynomial in one variable over the base ring, and $F'$ is its derivative.
\begin{enumerate}[label=(\roman*)]
\item $V$ contains all commutators $[x, y]$ with $x, y \in X$.
\item If $b \in X$ and $F'(b) = 0$, then $F(b) \in V$.
\item If $b \in X$, $F(b) \in V$ and $F'(b)$ is invertible, then $b \in V$.
\item $V$ contains all idempotents of $X$.
\item $V$ contains all elements $x^p$, where $x \in X$ and $p \geq 0$ is the characteristic of $X$.
\item If $a \in X$ has an inverse $a^{-1}$ in $X$ and $a^{-1} \in V$, then $a \in V$.
\item If $b \in X$ is invertible and $b^k \in V$ for some positive integer $k$, then $kb \in V$.
\end{enumerate}

\begin{theorem}[{\citep[Theorem 11]{Cheung2003}}]\label{thm:cheung}
Let $\mathfrak{A} = \mathrm{Tri}(A, M, B)$. If the following two conditions hold, then every Lie derivation of $\mathfrak{A}$ is proper.
\begin{enumerate}[label=(\Roman*)]
\item $M$ is a faithful left $A$-module and $Z(B) = \pi_B(Z(\mathfrak{A}))$; or $M$ is a faithful left $A$-module and $A = W(A)$; or every Lie derivation of $A$ is proper and $A = W(A)$.
\item $M$ is a faithful right $B$-module and $Z(A) = \pi_A(Z(\mathfrak{A}))$; or $M$ is a faithful right $B$-module and $B = W(B)$; or every Lie derivation of $B$ is proper and $B = W(B)$.
\end{enumerate}
\end{theorem}

The next corollary applies the theorem of Cheung to the $\sigma$-antisymmetrized derivations of $T$. % claim-ok: signpost for the corollary stated right after

\begin{corollary}\label{cor:triangular-lie}
Let $T = \mathrm{Tri}(A,M,B)$ satisfy conditions (I) and (II) of Theorem~\ref{thm:cheung}. Then every $\sigma$-antisymmetric derivation of $T$ for $\tau$ has the standard form $d + t$, and it is a derivation only if $t(1) = 0$.
\end{corollary}

\begin{proof}
By Proposition~\ref{prop:lie-classification}(i), the maps in question are the Lie derivations of $T$, which are proper by Theorem~\ref{thm:cheung}. The last claim is Proposition~\ref{prop:lie-classification}(iv).
\end{proof}

Corollary~\ref{cor:triangular-jordan} says that on such a triangular algebra the symmetrized product gives only derivations.
Example~\ref{ex:t2} shows that the antisymmetrized product can add a central term.
Take $T_2(K) = \mathrm{Tri}(K,K,K)$ over a field with $\operatorname{char} K \neq 2$.
Here $M = K$ is faithful on both sides.
Also $Z(T_2(K)) = KI$ projects onto $K$.
So Corollaries~\ref{cor:triangular-jordan} and~\ref{cor:triangular-lie} apply.

\begin{example}[Upper triangular $2 \times 2$ matrices]\label{ex:t2}
Let $T = T_2(\mathbb{F}_3)$.
Every additive map $T \to T$ is $\mathbb{F}_3$-linear.
So it is one of the $3^9$ maps given by a $3 \times 3$ matrix over $\mathbb{F}_3$.

An exhaustive computer check over all of them confirms Corollaries~\ref{cor:triangular-jordan} and~\ref{cor:triangular-lie}.
Every Jordan derivation is a derivation.
Every Lie derivation equals $\operatorname{ad}_x + t$ for some $x \in T$ and some central-valued $t$ that vanishes on commutators.

The term $t$ is needed.
Define $\theta$ by
\[
\theta\begin{pmatrix} a & m \\ 0 & b \end{pmatrix} = aI.
\]
It has central values and vanishes on commutators, whose diagonal is zero.
So it is a Lie derivation.
It is not a derivation, since
\[
\theta(e_1^2) = I \qquad \text{while} \qquad \theta(e_1)e_1 + e_1 \theta(e_1) = 2e_1.
\]
So on $T_2(\mathbb{F}_3)$, as on $M_n(K)$, we have $\JDer = \Der$, while $\LDer$ is strictly larger.
\end{example}

\begin{thebibliography}{9}

\bibitem[Adrabi and Bennis(2026)]{AdrabiBennis2026a}
Adrabi, A., Bennis, D.: Braided-symmetric derivations of algebras. Preprint (2026). \url{https://github.com/abderrahim-lectures/braided-sym-deriv}.

\bibitem[Alizadeh(2009)]{Alizadeh2009}
Alizadeh, R.: Jordan derivations of full matrix algebras. Linear Algebra Appl. \textbf{430}, 574--578 (2009).

\bibitem[Bre\v{s}ar(1988)]{Bresar1988}
Bre\v{s}ar, M.: Jordan derivations on semiprime rings. Proc. Amer. Math. Soc. \textbf{104}(4), 1003--1006 (1988).

\bibitem[Chambert-Loir(2021)]{ChambertLoir2021}
Chambert-Loir, A.: (Mostly) Commutative Algebra. Universitext. Springer, Cham (2021).

\bibitem[Cheung(2003)]{Cheung2003}
Cheung, W.-S.: Lie derivations of triangular algebras. Linear Multilinear Algebra \textbf{51}(3), 299--310 (2003).

\bibitem[Cusack(1975)]{Cusack1975}
Cusack, J.M.: Jordan derivations on rings. Proc. Amer. Math. Soc. \textbf{53}(2), 321--324 (1975).

\bibitem[Herstein(1957)]{Herstein1957}
Herstein, I.N.: Jordan derivations of prime rings. Proc. Amer. Math. Soc. \textbf{8}(6), 1104--1110 (1957).

\bibitem[Hochschild(1945)]{Hochschild1945}
Hochschild, G.: On the cohomology groups of an associative algebra. Ann. of Math. \textbf{46}(1), 58--67 (1945).

\bibitem[Jacobson(1937)]{Jacobson1937}
Jacobson, N.: Abstract derivation and Lie algebras. Trans. Amer. Math. Soc. \textbf{42}(2), 206--224 (1937).

\bibitem[Lam(2001)]{Lam2001}
Lam, T.Y.: A First Course in Noncommutative Rings, 2nd edn. Graduate Texts in Mathematics 131. Springer, New York (2001).

\bibitem[Martindale(1964)]{Martindale1964}
Martindale, W.S., III: Lie derivations of primitive rings. Michigan Math. J. \textbf{11}(2), 183--187 (1964).

\bibitem[Rotman(2009)]{Rotman2009}
Rotman, J.J.: An Introduction to Homological Algebra, 2nd edn. Universitext, Springer, New York (2009).

\bibitem[Zhang and Yu(2006)]{ZhangYu2006}
Zhang, J.-H., Yu, W.-Y.: Jordan derivations of triangular algebras. Linear Algebra Appl. \textbf{419}(1), 251--255 (2006).

\bibitem[Wang(2019)]{Wang2019}
Wang, Y.: Lie (Jordan) derivations of arbitrary triangular algebras. Aequationes Math. \textbf{93}, 1221--1229 (2019).

\bibitem[Weibel(1994)]{Weibel1994}
Weibel, C.A.: An Introduction to Homological Algebra. Cambridge Stud. Adv. Math. 38, Cambridge University Press, Cambridge (1994).

\end{thebibliography}

\end{document}
